% =============================================================================
%  第 2 章  行列式
% =============================================================================
\chapter{行列式}\label{chap:det}

\epigraph{纯数学的对象是现实世界的空间形式和数量关系。}{恩格斯《反杜林论》}
\chaptertoc

上一章用消元法回答了线性方程组的可解性问题。对于方程个数与未知数个数相同的方程组，还有
一个更紧凑的判据：方程组有唯一解，当且仅当一个由系数算出的数——行列式——不等于零。本章
从平面上的面积出发引入行列式，讨论它的性质与计算，并给出克拉默法则。

\section{面积与有向面积}\label{sec:area}

设平面上的两个向量 $\vb{u}=(a_1,a_2)$，$\vb{v}=(b_1,b_2)$ 不共线，它们张成一个平行四边形。
如\cref{fig:area}(a)，用外接矩形的面积减去四周的两对直角三角形和两个小矩形，得到
\[
  (a_1+b_1)(a_2+b_2)-a_1a_2-b_1b_2-2a_2b_1=a_1b_2-a_2b_1 .
\]
这里假定了 $\vb{v}$ 位于 $\vb{u}$ 的逆时针一侧；若交换 $\vb{u}$ 与 $\vb{v}$ 的位置，同样的式子给出
面积的相反数。于是我们定义 $\vb{u}$，$\vb{v}$ 张成的\term[you xiang mian ji]{有向面积}为
\begin{equation}\label{eq:2x2}
  D(\vb{u},\vb{v})=\begin{vmatrix}a_1&a_2\\b_1&b_2\end{vmatrix}=a_1b_2-a_2b_1 ,
\end{equation}
它的绝对值是平行四边形的面积，符号则记录了从 $\vb{u}$ 到 $\vb{v}$ 的转向。

\begin{figure}[htbp]
\centering
\begin{subfigure}{.48\textwidth}
\centering
\begin{tikzpicture}[scale=.8]
  \coordinate (O) at (0,0); \coordinate (U) at (3,1);
  \coordinate (V) at (1,2.5); \coordinate (W) at (4,3.5);
  \draw[mj grid] (0,0) grid (4,3.5);
  \fill[mj fill] (O) -- (U) -- (W) -- (V) -- cycle;
  \draw[mj thin] (0,0) rectangle (4,3.5);
  \draw[mj aux] (3,0) -- (U) -- (4,1) (1,3.5) -- (V) -- (0,2.5);
  \draw[mj curve, ->] (O) -- (U) node[midway, below right=-2pt] {$\vb{u}$};
  \draw[mj curve2, ->] (O) -- (V) node[midway, above left=-2pt] {$\vb{v}$};
  \draw[mjAccent, line width=.6pt] (U) -- (W) -- (V);
  \node[below] at (1.5,0) {\footnotesize$a_1$};
  \node[below] at (3.5,0) {\footnotesize$b_1$};
  \node[left] at (0,1.25) {\footnotesize$b_2$};
  \node[left] at (0,3) {\footnotesize$a_2$};
\end{tikzpicture}
\caption{面积 $=a_1b_2-a_2b_1$}
\end{subfigure}\hfill
\begin{subfigure}{.48\textwidth}
\centering
\begin{tikzpicture}[scale=.8]
  \coordinate (O) at (0,0); \coordinate (U) at (3,0);
  \coordinate (V) at (.6,2.2); \coordinate (V2) at (2.4,2.2);
  \fill[mj fill] (O) -- (U) -- ($(U)+(V)$) -- (V) -- cycle;
  \fill[mj fill2, opacity=.9] (O) -- (U) -- ($(U)+(V2)$) -- (V2) -- cycle;
  \draw[mj aux] (-.4,2.2) -- (5.8,2.2);
  \draw[mj curve, ->] (O) -- (U) node[midway, below] {$\vb{u}$};
  \draw[mj curve2, ->] (O) -- (V) node[midway, left] {$\vb{v}$};
  \draw[mj curve2, ->] (O) -- (V2) node[pos=.62, right=1pt] {$\vb{v}+t\vb{u}$};
  \draw[mjAccent, line width=.5pt] (U) -- ($(U)+(V)$) -- (V);
  \draw[mjSecond, line width=.5pt] (U) -- ($(U)+(V2)$) -- (V2);
\end{tikzpicture}
\caption{同底等高：$D(\vb{u},\vb{v}+t\vb{u})=D(\vb{u},\vb{v})$}
\end{subfigure}
\caption{平行四边形的有向面积}
\label{fig:area}
\end{figure}

有向面积有三条基本性质，它们都可以从 \cref{eq:2x2} 直接验证，也都有明显的几何意义：
\begin{enumerate}
  \item \textbf{多重线性}：$D$ 对每个变元都是线性的，例如
        $D(c\vb{u}+c'\vb{u}',\vb{v})=cD(\vb{u},\vb{v})+c'D(\vb{u}',\vb{v})$；
  \item \textbf{交错性}：$D(\vb{u},\vb{u})=0$，即退化的平行四边形面积为零；
  \item \textbf{规范性}：$D(\vb{e}_1,\vb{e}_2)=1$，即单位正方形的面积为 $1$。
\end{enumerate}
由 (1)(2) 立即得到 $D(\vb{u},\vb{v}+t\vb{u})=D(\vb{u},\vb{v})$，这正是\cref{fig:area}(b) 所示的
“同底等高的平行四边形面积相等”。下一节我们看到，这三条性质完全刻画了行列式。

\section{\texorpdfstring{$n$}{n} 阶行列式}\label{sec:det}

\begin{definition}\label{def:det}
设 $\det$ 是定义在全体 $n$ 阶方阵上的函数，把它看作 $n$ 个行向量的函数。若它满足
\begin{enumerate}
  \item 对每一行都是线性的（\term[duo chong xian xing]{多重线性}）；
  \item 若有两行相同，则函数值为零（\term[jiao cuo xing]{交错性}）；
  \item $\det I_n=1$，
\end{enumerate}
则称 $\det A$ 为方阵 $A$ 的\term[hang lie shi]{行列式}，也记作 $\abs{A}$。
\end{definition}

\begin{theorem}\label{thm:leibniz}
满足\cref{def:det} 的函数存在且唯一，它由下面的展开式给出：
\begin{equation}\label{eq:leibniz}
  \det A=\sum_{\sigma\in S_n}\sgn(\sigma)\,a_{1\sigma(1)}a_{2\sigma(2)}\dotsm a_{n\sigma(n)},
\end{equation}
其中 $S_n$ 是 $\{1,\dotsc,n\}$ 的全体置换；$\sigma$ 为偶置换时 $\sgn\sigma=1$，为奇置换时 $\sgn\sigma=-1$。
\end{theorem}

\begin{proof}[证明概要]
先看唯一性。由交错性与线性，交换两行使行列式变号：
$0=\det(\dotsc,\vb{a}+\vb{b},\dotsc,\vb{a}+\vb{b},\dotsc)=\det(\dotsc,\vb{a},\dotsc,\vb{b},\dotsc)+\det(\dotsc,\vb{b},\dotsc,\vb{a},\dotsc)$。
把第 $i$ 行写成 $\sum_ja_{ij}\vb{e}_j$，利用多重线性完全展开，得到 $n^n$ 项
$a_{1j_1}\dotsm a_{nj_n}\det(\vb{e}_{j_1},\dotsc,\vb{e}_{j_n})$。下标有重复的项由交错性为零；
其余各项对应置换 $\sigma$，经过若干次对换化为 $\det I_n$，得到 $\sgn\sigma$。这就证明了
函数必为 \cref{eq:leibniz}。反过来，直接验证 \cref{eq:leibniz} 满足三条公理即得存在性。
\end{proof}

当 $n=2$ 时，\cref{eq:leibniz} 就是 \cref{eq:2x2}。当 $n=3$ 时，它包含 $3!=6$ 项：
\[
\begin{vmatrix}a_{11}&a_{12}&a_{13}\\a_{21}&a_{22}&a_{23}\\a_{31}&a_{32}&a_{33}\end{vmatrix}
=a_{11}a_{22}a_{33}+a_{12}a_{23}a_{31}+a_{13}a_{21}a_{32}
-a_{13}a_{22}a_{31}-a_{11}a_{23}a_{32}-a_{12}a_{21}a_{33},
\]
可以按\cref{fig:sarrus} 所示的\term[dui jiao xian fa ze]{对角线法则}记忆。注意这一法则只适用于三阶行列式。

\begin{figure}[htbp]
\centering
\begin{tikzpicture}[x=11mm, y=8mm]
  \foreach \i in {1,2,3} \foreach \j in {1,2,3,4,5} {
    \pgfmathtruncatemacro{\jj}{mod(\j-1,3)+1}
    \ifnum\j>3
      \node[mjGray] (a\i\j) at (\j,-\i) {$a_{\i\jj}$};
    \else
      \node (a\i\j) at (\j,-\i) {$a_{\i\jj}$};
    \fi
  }
  \draw[mjInk, line width=.6pt] (.55,-.6) -- (.55,-3.4) (3.45,-.6) -- (3.45,-3.4);
  \foreach \j [evaluate=\j as \k using int(\j+2)] in {1,2,3}
    \draw[mjAccent, line width=.7pt, opacity=.75, shorten <=-3pt, shorten >=-3pt]
      (a1\j) -- (a3\k);
  \foreach \j [evaluate=\j as \k using int(\j-2)] in {3,4,5}
    \draw[mjSecond, line width=.7pt, opacity=.75, dashed, shorten <=-3pt, shorten >=-3pt]
      (a1\j) -- (a3\k);
  \node[mjAccent, font=\small] at (6.3,-1.4) {实线：取正号};
  \node[mjSecond, font=\small] at (6.3,-2.6) {虚线：取负号};
\end{tikzpicture}
\caption{三阶行列式的对角线法则：把前两列抄在右侧，沿对角线相乘}
\label{fig:sarrus}
\end{figure}

直接用 \cref{eq:leibniz} 计算高阶行列式是不现实的（$n!$ 增长极快）。实际计算依靠下面的性质，
把行列式化为三角形。

\begin{theorem}\label{thm:det-properties}
设 $A$，$B$ 是 $n$ 阶方阵，则
\begin{enumerate}
  \item $\det A\T=\det A$，因而下面关于行的结论对列同样成立；
  \item 对换两行，行列式变号；某行乘以 $c$，行列式乘以 $c$；倍加变换不改变行列式；
  \item 三角形矩阵的行列式等于主对角线上元素之积；
  \item $\det(AB)=\det A\cdot\det B$。
\end{enumerate}
\end{theorem}

性质 (2) 由定义直接得到；(1)(3)(4) 的证明见\parencite[第 2 章]{pku2013}。

\begin{example}\label{ex:det4}
计算 $D=\begin{vmatrix}1&2&3&4\\2&3&4&1\\3&4&1&2\\4&1&2&3\end{vmatrix}$。
\end{example}

\begin{solution}
各行元素之和都是 $10$。把第 2、3、4 列都加到第 1 列，再提出公因子 $10$：
\begin{align*}
D&=10\begin{vmatrix}1&2&3&4\\1&3&4&1\\1&4&1&2\\1&1&2&3\end{vmatrix}
\xlongequal[r_4-r_1]{r_2-r_1,\ r_3-r_1}
10\begin{vmatrix}1&2&3&4\\0&1&1&-3\\0&2&-2&-2\\0&-1&-1&-1\end{vmatrix}\\
&\xlongequal[r_4+r_2]{r_3-2r_2}
10\begin{vmatrix}1&2&3&4\\0&1&1&-3\\0&0&-4&4\\0&0&0&-4\end{vmatrix}
=10\times1\times1\times(-4)\times(-4)=160 .
\end{align*}
\end{solution}

\section{克拉默法则}\label{sec:cramer}

\begin{theorem}[克拉默法则]\label{thm:cramer}
设 $A$ 是 $n$ 阶方阵，$\det A\ne0$，则线性方程组 $A\vb{x}=\vb{b}$ 有唯一解
\begin{equation}\label{eq:cramer}
  x_j=\frac{\det A_j}{\det A},\qquad j=1,2,\dotsc,n,
\end{equation}
其中 $A_j$ 是把 $A$ 的第 $j$ 列换成 $\vb{b}$ 所得的矩阵。
\end{theorem}

\begin{proof}
由\cref{thm:det-properties}，初等行变换只把行列式乘以非零数，所以 $A$ 的简化行阶梯形 $R$
满足 $\det R\ne0$。$R$ 若有零行则行列式为零，故 $R=I_n$，没有自由变量，由\cref{thm:existence}
方程组有唯一解 $\vb{x}$。记 $A$ 的列为 $\vb{a}_1,\dotsc,\vb{a}_n$，则 $\vb{b}=\sum_kx_k\vb{a}_k$。
由行列式对列的多重线性与交错性，
\[
  \det A_j=\det\Bigl(\vb{a}_1,\dotsc,\sum_kx_k\vb{a}_k,\dotsc,\vb{a}_n\Bigr)
  =\sum_kx_k\det(\vb{a}_1,\dotsc,\vb{a}_k,\dotsc,\vb{a}_n)=x_j\det A ,
\]
因为 $k\ne j$ 的各项中有两列相同。
\end{proof}

\begin{example}\label{ex:he-cramer}
用克拉默法则重解\cref{sec:jiuzhang} 中的方程组 \cref{eq:he}。
\end{example}

\begin{solution}
计算四个三阶行列式：
\begin{alignat*}{2}
\det A&=\begin{vmatrix}3&2&1\\2&3&1\\1&2&3\end{vmatrix}=12, &\qquad
\det A_1&=\begin{vmatrix}39&2&1\\34&3&1\\26&2&3\end{vmatrix}=111,\\
\det A_2&=\begin{vmatrix}3&39&1\\2&34&1\\1&26&3\end{vmatrix}=51, &\qquad
\det A_3&=\begin{vmatrix}3&2&39\\2&3&34\\1&2&26\end{vmatrix}=33 .
\end{alignat*}
于是 $x=\frac{111}{12}=\frac{37}{4}$，$y=\frac{51}{12}=\frac{17}{4}$，$z=\frac{33}{12}=\frac{11}{4}$，与消元法的结果一致。
\end{solution}

\begin{remark}
克拉默法则在理论上很重要：它表明解的每个分量都是系数的有理函数。但用它来计算并不经济：
即便每个行列式都用消元法计算，也要计算 $n+1$ 个行列式，运算量约为高斯消元法的 $n$ 倍。
\end{remark}

\section{范德蒙德行列式}\label{sec:vandermonde}

\begin{theorem}\label{thm:vandermonde}
对任意数 $x_1,\dotsc,x_n$（$n\ge2$），有
\begin{equation}\label{eq:vandermonde}
  V_n=\begin{vmatrix}
  1&x_1&x_1^2&\cdots&x_1^{n-1}\\
  1&x_2&x_2^2&\cdots&x_2^{n-1}\\
  \vdots&\vdots&\vdots&&\vdots\\
  1&x_n&x_n^2&\cdots&x_n^{n-1}
  \end{vmatrix}
  =\prod_{1\le i<j\le n}(x_j-x_i).
\end{equation}
\end{theorem}

\begin{proof}
对 $n$ 归纳。$n=2$ 时显然。设结论对 $n-1$ 成立。依次把第 $n-1,n-2,\dotsc,1$ 列的 $x_1$ 倍从
第 $n,n-1,\dotsc,2$ 列中减去，第一行变为 $(1,0,\dotsc,0)$，第 $i$ 行（$i\ge2$）变为
$\bigl(1,\;x_i-x_1,\;x_i(x_i-x_1),\;\dotsc,\;x_i^{n-2}(x_i-x_1)\bigr)$。按第一行展开，再从第 $i$ 行
提出公因子 $x_i-x_1$，得
\[
  V_n=\prod_{i=2}^n(x_i-x_1)\cdot V_{n-1}(x_2,\dotsc,x_n),
\]
由归纳假设即得 \cref{eq:vandermonde}。
\end{proof}

\begin{corollary}[插值多项式的存在唯一性]\label{cor:interpolation}
设 $x_1,\dotsc,x_n$ 互不相同，则对任意 $y_1,\dotsc,y_n$，存在唯一的次数不超过 $n-1$ 的多项式
$p$，使 $p(x_i)=y_i$（$i=1,\dotsc,n$）。
\end{corollary}

\begin{proof}
设 $p(x)=c_0+c_1x+\dotsb+c_{n-1}x^{n-1}$，条件 $p(x_i)=y_i$ 是关于 $c_0,\dotsc,c_{n-1}$ 的线性方程组，
其系数矩阵的行列式正是 $V_n\ne0$。由克拉默法则，它有唯一解。
\end{proof}

例如，过四点 $(0,1)$，$(1,3)$，$(2,2)$，$(3,5)$ 的\term[cha zhi duo xiang shi]{插值多项式}是
$p(x)=1+\frac{35}{6}x-5x^2+\frac{7}{6}x^3$（\cref{fig:interp}）。

\begin{figure}[htbp]
\centering
\begin{tikzpicture}
\begin{axis}[mojing, width=.62\textwidth, height=.38\textwidth,
    xmin=-.5, xmax=3.6, ymin=-.4, ymax=6.2, xtick={1,2,3}, ytick={1,2,3,4,5},
    xlabel=$x$, ylabel=$y$, samples=100, domain=-.3:3.3]
  \addplot[mjAccent] {1+35/6*x-5*x^2+7/6*x^3};
  \addplot[only marks, mark=*, mark size=1.8pt, mjInk] coordinates {(0,1) (1,3) (2,2) (3,5)};
\end{axis}
\end{tikzpicture}
\caption{过四个点的唯一一条三次插值曲线}
\label{fig:interp}
\end{figure}

\begin{exercises}
\exgroup{A 组}
\begin{enumerate}
  \item 计算 $\begin{vmatrix}2&1&3\\1&-1&2\\3&2&1\end{vmatrix}$。
  \item 计算 $n=4$ 时的行列式 $\begin{vmatrix}a&b&b&b\\b&a&b&b\\b&b&a&b\\b&b&b&a\end{vmatrix}$。
  \item 用克拉默法则解方程组 $\left\{\begin{aligned}2x+y&=5\\x-3y&=-1\end{aligned}\right.$。
  \item 不用\cref{thm:vandermonde}，直接证明
        $\begin{vmatrix}1&1&1\\a&b&c\\a^2&b^2&c^2\end{vmatrix}=(b-a)(c-a)(c-b)$。
\end{enumerate}
\exgroup{B 组}
\begin{enumerate}
  \item 设 $A$ 是奇数阶反对称矩阵，即 $A\T=-A$。证明 $\det A=0$。
  \item 证明：平面上三点 $(x_i,y_i)$（$i=1,2,3$）共线的充要条件是
        $\begin{vmatrix}x_1&y_1&1\\x_2&y_2&1\\x_3&y_3&1\end{vmatrix}=0$。
  \item 设 $A$ 是元素都为整数的 $n$ 阶方阵。证明：对每个整数向量 $\vb{b}$，方程组 $A\vb{x}=\vb{b}$
        都有整数解，当且仅当 $\det A=\pm1$。
\end{enumerate}
\end{exercises}

\begin{history}
行列式的概念几乎同时出现在东西方。日本数学家关孝和在 1683 年的著作中已经给出了相当于
行列式展开的计算方法；十年后的 1693 年，莱布尼茨在致洛必达的信中也用到了同样的思想。
1750 年，克拉默在研究代数曲线时发表了以他命名的法则\parencite{kline1972}。范德蒙德于 1772 年
第一个把行列式作为独立的对象来研究——有趣的是，以他命名的那个行列式并没有出现在他的著作中。
1812 年，柯西系统地建立了行列式理论，证明了乘法公式 $\det(AB)=\det A\det B$，并在今天的意义下
使用了“行列式”（déterminant）这一名称；在此之前，高斯曾用这个词指二次型的判别式。
\end{history}
